\newglossaryentry{bnchoice}{name={$\bnchoice$}, description={non-deterministic choice}}

\chapter{Pattern Language}\label{ch:language}
\minitoc

With the detail of the graph model for patterns and layouts fully discussed.
We give a language extension for pattern matching to general {\OO} languages.
We then present the type-checking rules and operational semantics, and prove the type safety to reason about the soundness of our model.

\section{Syntax}

\begin{table}[t]
\caption{Abstract Syntax of Pattern Matching}\label{tab:syntax}
\centering
\[
\begin{array}[t]{ccl@{\quad}l}
  P &\cceq& [r:] \; (K \orr r \orr V) & \mbox{pattern} \\
  K &\cceq& T \; \overline{a = P} & \mbox{class pattern} \\
  r & & & \mbox{pattern reference} \\
  V & & & \mbox{value set} \\
  T & & & \mbox{class type} \\
  a & & & \mbox{attribute} \\
  J &\cceq& P \orr \bgconj\overline{P} \orr \bgdisj\overline{P} & \mbox{pattern junction} \\
  M &\cceq& e \; \ltimes \; \overline{J \rightarrow S} & \mbox{match statement} \\
  S & & & \mbox{statement}
\end{array}
\]
\end{table}

The abstract syntax to define patterns and matches is given in \autoref{tab:syntax}.
The syntax is much like conventional ones found in other OO and functional languages, with the addition of pattern references.
The pattern references play two roles, to be referred to from inside the pattern to produce aliases and cycles, and
to be referred to from the match-selected statement as object bindings --- variables bound to components of matching objects.

A pattern $P$ can have an optional label, which is a pattern reference $r$,
followed by a pattern specification in one of the three forms ---
a class pattern, a pattern reference or a set $V$ of values.
A class pattern $K$ is a class type $T$ followed by zero or more attribute specifications,
each of which is an attribute name $a$ equal to an inner pattern $P$.
For example, the pattern $p_4$ shown in \autoref{fig:junction} can be expressed as
\[
T_1(b = T_2, c = y: T_3, d = y).
\]
The more complex pattern shown in \autoref{fig:obj-pattern} is written as
\[
T(l = w: T(l = x: T, r = y: T(l = w, r = z)), r = z: T).
\]
A match statement evaluates an expression $e$ and matches the result against many cases in turn.
Each case consists of a pattern junction $J$, which can be a single pattern,
a pattern conjunction $\bgconj\overline{P}$ or a pattern disjunction $\bgdisj\overline{P}$,
and an associated statement $S$.
The first successfully matched case selects the associated statement to execute,
with the new object bindings introduced by the references declared in the matched pattern junction.

\section{Pattern Graph Construction}

To carry out the operations of pattern matching and check for its type safety based on the graph model described in \autoref{ch:graph},
we must first construct the pattern graph for each pattern specified in the abstract syntax.
The construction is straightforward following the structure of the specification, with three components to output ---
a layout graph $G$, a map $T$ from nodes to types and a map $R$ from pattern references to nodes.
\autoref{algo:graph-cons} describes the procedure in detail,
where, for clarity, we enclose syntax elements in guillemets $\syn{}$ to distinguish them from other variables.

\begin{algorithm}[t]
\caption{Construction of Pattern Graph}\label{algo:graph-cons}
\DontPrintSemicolon
\SetKwFunction{Parse}{parse}
\KwIn{pattern:~$\syn{P}$.}
\KwOut{layout graph:~$G$, type~map:~$T$ and reference~map:~$R$.}
\KwEx{$\SynUndefRef,\SynRedefRef$.}
\Begin{
  $N,E,T,R \gets \emptyset$\;
  $p \gets {}$ \Parse{$\syn{P}$}\;
  \lIRaise{$\SynUndefRef$}{$\ran(R) \nsubseteq N$}
  $G \gets \vect{N,E,p}$
}
\medskip
\Fn{\Parse{$\syn{P}$} ${} \to p$}{
  \uIf{$\syn{P \redto r:Q}$}{
    $p \gets {}$ \Parse{$\syn{Q}$}\;
    \If{$\syn{r} \in \dom(R)$}{
      $r \gets R(\syn{r})$\;
      \lIRaise{$\SynRedefRef$}{$r \in N$}
      $E \gets E[p/r], R \gets R[p/r]$\;
    }
    $R \gets R \uplus \sett{\syn{r} \mapsto p}$
  }
  \uElseIf{$\syn{P \redto T \; \overline{a = Q}}$}{
    $p \gets \text{new node}$\;
    $N \gets N \cup \sett{p}$\;
    $T \gets T \cup \sett{p \mapsto \syn{T}}$\;
    $\overline{q \gets \mbox{\Parse{$\syn{Q}$}}}$\;
    $E \gets E \cup \sett{\overline{p\la{\syn{a}}q}}$
  }
  \uElseIf{$\syn{P \redto r}$}{
    \eIf{$\syn{r} \in \dom(R)$}{
      $p \gets R(\syn{r})$
    }{
      $p \gets \text{new node}$\;
      $R \gets R \cup \sett{\syn{r} \mapsto p}$
    }
  }
  \ElseIf{$\syn{P \redto V}$}{
    $p \gets \text{new node}$\;
    $N \gets N \cup \sett{p}$\;
    $T \gets T \cup \sett{p \mapsto \syn{V}}$
  }
}
\end{algorithm}

In the algorithm, function {\Parse} recursively decomposes a pattern specification $\syn{P}$ into a graph with root $p$
by updating the nodes $N$, the edges $E$, the type map $T$ and the reference map $R$.
The {\Parse} function performs case analysis on the pattern specification.
First, when it is a reference definition $\syn{r:Q}$, we recursively obtain the root $p$ of $\syn{Q}$, and associate it with the reference $\syn{r}$ in $R$.
It is possible for a reference to be used before definition, we generate a placeholder node $r$ for the forward use, and replace it with the real node at this definition stage.
A placeholder is not put into the node set $N$, therefore we can check if $r$ is indeed a placeholder, or a node that has defined the reference already.
For a placeholder $r$, we need to substitute the real node $p$ for it in the edges $E$, for all the source and target nodes equal to $r$,
and also in the reference map $R$, for all the image nodes equal to $r$. Note that we allow a reference to label another reference.
Second, when it is a class pattern $\syn{T\;\overline{a=Q}}$, we create a new node $p$ and associate it with type $\syn{T}$.
We add $p$ to the node set $N$ and link the outgoing edges to the roots recursively obtained from $\overline{\syn{Q}}$, with labels equal to the attribute names $\overline{\syn{a}}$.
Third, when it is a reference occurrence $\syn{r}$, we try to look up $R$ for its corresponding node, or create a new placeholder without placing it into the node set $N$.
Last, when it is a value set $\syn{V}$, we simply create a new node and associate it with $\syn{V}$ in the type map $T$.

The graph construction merely creates nodes and links between the nodes. This process is always successful.
However, with the addition of references, there are two exceptional cases ---
a reference is defined more than once ($\SynRedefRef$), and a reference occurs but not defined ($\SynUndefRef$).
The former is detected at the time of parsing a label prefix, while the latter is detected when the parsing is complete,
by checking whether all the placeholders have been resolved.

\section{Type-Checking}

\begin{table}[t]
\caption{Type-Checking Rules}\label{tab:tc}
\centering
\normalsize
\begin{mathpar}
\inferrulel[(Node)]
  {\Env \vdash t = T(p) \\ \overline{p\la{a}q \in E} \\ p:t,\Env \xt E,T \vdash \overline{q:s} \\
   \Env \vdash \overline{t.a:s' \wedge s \subty s'}\Olf{2}}
  {\Env \xt E,T \vdash p:t}
\and
\inferrulel[(Graph)]
  {P = \vect{N,E,p,T} \\ \Env \xt E,T \vdash p:t}
  {\Env \vdash P:t}
\and
\inferrulel[(Ref)]
  {R(r) = p \\ P = \vect{N,E,p,T} \\ \Env; E,T \vdash p:t}
  {\Env \xt P,R \vdash r:t}
\and
\inferrulel[(Pattern)]
  {\mbox{\autoref{algo:graph-cons}}: \syn{P} \to G,T,R \\
   P = \vect{G,T} \\ \Env \vdash P:t \\
   \overline{r} = \dom(R) \\ \Env \xt P,R \vdash \overline{r:s}}
  {\Env \vdash \syn{P}:t,\overline{r:s} \xt P,R}
\and
\inferrulel[(Conj)]
  {\Env \vdash \overline{\syn{P}:t, \overline{r:s} \xt P,R} \\
   \mbox{\autoref{algo:conj-dfs-cons}}:\Env,\overline{P} \to \overline{f},\co{T} \\
   \overline{s' = \min(s,\co{T}(f(R(r))))} \\ \hat{t} = \co{T}(f_1(R_1(\Root(P_1)))) \\ \overline{t' = \min(t,\hat{t})}}
  {\Env \vdash \syn{\bgconj\overline{P}}:\hat{t},\overline{\syn{P}:t'},\overline{r:s'} \xt \overline{P},\overline{f},\overline{R}}
\and
\inferrulel[(Disj)]
  {\Env \vdash \overline{\syn{P}:t,\overline{r:s} \xt P,R} \\
   \mbox{\autoref{algo:disj-dfs-cons}}:\Env,\overline{P} \to \overline{f},\ov{T} \\
   \overline{s' = \ov{T}(f(R(r)))} \\ \hat{t} = \ov{T}(f_1(R_1(\Root(P_1)))}
  {\Env \vdash \syn{\bgdisj\overline{P}}:\hat{t},\overline{\syn{P}:t},\overline{r:s'} \xt \overline{P},\overline{f},\overline{R}}
\and
\inferrulel[(Case)]
  {\Env \vdash \syn{J}:t,\overline{r:s} \xt \overline{P},\overline{f},\overline{R} \\
   \overline{r:s},\Env \vdash \syn{S}}
  {\Env \vdash \syn{J \to S}: t}
\and
\inferrulel[(Match)]
  {\Env \vdash \overline{\syn{J \to S}: t} \\
   \Env \vdash e: t' \\ \Env \vdash \overline{t' \subty t \vee t \subty t'}}
  {\Env \vdash \syn{e \ltimes \overline{J \to S}}}
\end{mathpar}
\end{table}

The pattern matching extension is type-checked in two parts, the pattern and the match.
A pattern can be a class pattern, a reference, a pattern conjunction or a pattern disjunction.
A class pattern is converted to a pattern graph and checked by recursively reducing to each node and the outgoing edges.
A node in a pattern graph is well-typed when all its outgoing edges are labeled with the attributes of the associated class,
and also the targets are all well-typed.
To avoid infinite loops in the {\DFS}, we add the visited nodes to the typing environment before going into recursion.
A reference is checked by resolving it to a node with the reference map.
A class pattern is well-typed only when the corresponding pattern graph and all the references declared in the pattern are well-typed.

The type-checking rules are given in \autoref{tab:tc}.
Each typing rule involves a typing environment $\Env$ and optionally some additional structures, separated by a semicolon, as the context.
The result of a rule is the well-typedness and type-assignment of the checked element, possibly with some extra output, also separated by a semicolon, during the checking.
This output can be necessary in the subsequent procedures.

For a pattern conjunction, we check the individual patterns first, and try to construct the match conjunction.
Then, we promote the type of each referenced node to the infimum type associated with the union.
If such infimum type does not exist, we fall back to the original type of the node.
For a pattern disjunction, we try to construct the match disjunction.
Only a referenced node that occurs commonly in all the component patterns has a type, which must be the supremum type associated with the intersection.
If such supremum type does not exist, the referenced node has $\NoType$, and subsequent uses of the reference will be ill-typed.
This includes the case where a referenced node does not occur in all the component patterns.

A match is a test of an expression against multiple cases.
Therefore, to type-check a match, we need to check the individual cases, and whether the expression is type-compatible with all the cases.
Each case consists of a pattern and a statement.
The checking of the pattern produces the type of the pattern and a series of type-assigned references, as the additional variables available to the statement.
We check the type of the pattern against the expression to match, and check the statement under the context with those additional variables declared.

\section{Operational Semantics}

\begin{table}[t]
\caption{State Transition Rules}\label{tab:st}
\centering
\normalsize
\begin{mathpar}
\inferrulel[(Match)]
  {\eval(\syn{e}) = p \\ P = \Sta \rrt p}
  {\vect{\syn{e \ltimes \overline{J \to S}}, \Sta} \to \vect{P \ltimes \syn{\overline{J \to S}}, \Sta}}
\and
\inferrulel[(Match-None)]
  {}
  {\vect{P \ltimes \syn{\emptyset}, \Sta} \to \Sta}
\and
\inferrulel[(Match-Some)]
  {\syn{\overline{J \to S}} = \syn{J_1 \to S_1, \overline{J' \to S'}} \\ \Env \vdash \syn{J_1} \xt \overline{P_1},\overline{f_1},\overline{R_1}}
  {\inferrule
    {\mb(P,\syn{J_1},\overline{P_1},\overline{f_1},\overline{R_1}) = (\overline{r},\overline{q})}
    {\vect{P \ltimes \syn{\overline{J \to S}}, \Sta} \to \vect{\syn{\Var(\overline{r=q});S;\EndV(\overline{r})}, \Sta}} \\
   \inferrule
    {\mb(P,\syn{J_1},\overline{P_1},\overline{f_1},\overline{R_1}) = \emptyset}
    {\vect{P \ltimes \syn{\overline{J \to S}}, \Sta} \to \vect{P \ltimes \syn{\overline{J' \to S'}}, \Sta}}}
\end{mathpar}

\begin{mathpar}
\mb(P,\syn{J_1},\overline{P_1},\overline{f_1},\overline{R_1}) \qdef
\hfill\and
\inferrulel[(Match-One)]
  {\syn{J_1} = \syn{P} \\ \Env \vdash P \gmat[f] P_1}
  {\overline{r} = \dom(R_1) \\ \overline{q = f(R_1(r))}}
\and
\inferrulel[(Match-Conj)]
  {\syn{J_1} = \syn{\bgconj\overline{P}} \\ \Env \vdash \overline{P \gmat[f] P_1} \\
   R' = \bgunion\overline{R_1} \\ f' = \bgunion\overline{f}}
  {\overline{r} = \dom(R') \\ \overline{q = f'(R'(r))}}
\and
\inferrule[(Match-Disj)]
  {\syn{J_1} = \syn{\bgdisj\overline{P}} \\ \Env \vdash \exists P' \in \overline{P_1} \suchthat P \gmat[f] P' \\
   R' = \bgunion\overline{R_1} \\ f_1' = \bgunion\overline{f_1} \\ \overline{\ov{u} = f_1'(R'(r))}\Olf{2}}
  {\overline{r} = \dom(R') \\ \overline{q = \mgls{bnchoice}\sett{f(u) \mid f_1'(u) = \ov{u}}}}
\and
\inferrulel[(Mismatch)]
  {\mbox{otherwise}}
  {\emptyset}
\end{mathpar}
\end{table}

The operation of a pattern match statement is to evaluate the expression,
then to match the result against the list of cases in turn,
and finally to select the statement in the matched case to execute.
In the graph model, the expression $\syn{e}$ to match is evaluated to a node in the state graph,
and then the pattern graph $P$ rooted at this node is extracted from the state graph.
The pattern graph is matched in turn against all the cases until a matched case is found,
otherwise there is no statement to execute in this match.

The operational semantics is given by the state transition rules listed in \autoref{tab:st}.
Before we can match the pattern graph of the expression,
we need to obtained the pattern graphs of the patterns by taking the output of the type-checking procedures.
Beside to match by constructing the match morphisms, we also need to obtain the references declared in the patterns, as well as their corresponding nodes.
This task is done in the $\mb$ function, divided into three cases, a single pattern, a conjunction, or a disjunction.
For a single pattern, a reference $r$ is mapped to a node of the pattern graph $P_1$ by the reference map $R_1$,
then we can use the match morphism $f$ from the pattern graph to the expression graph $P$ to obtain the corresponding node in the state graph.
For a conjunction, since the match of the expression graph must be done against all the component pattern graphs,
we have a series of match morphisms from different pattern graphs to the same expression graph, their domains do not intersect with each other.
We can simply union these morphisms into one to obtain the nodes of the references.

A disjunction is more complicated, because there requires only one match of the component pattern graphs,
and a node in the intersection is possibly de-aliased to two or more nodes in a component pattern.
We have restricted the references in a disjunction only to those nodes in the intersection, by assigning $\NoType$ to the other nodes.
Therefore, we are left with a one-to-many mapping from the intersection to the matched pattern, and thus to the expression graph.
We use the non-deterministic choice $\nchoice$ to bind a reference to one of the possible corresponding nodes in the state graph.
Note that, in function $\mb$, $f_1'$ is the union of the morphisms from the patterns to their intersection.
Once we have the bindings of the references to the nodes in the state graph,
we can declare the additional variables as a normal variable declaration, $\Var\ldots\EndV$, and execute the statement of the matched case in this scope.

\section{Type Safety}

The type safety of pattern matching means a well-typed pattern match statement cannot go wrong at runtime.
Since the only new operation introduced by pattern matching is the new bindings of pattern references to runtime objects,
the statement of the matched case is type-safe when the runtime objects truly have the types associated with the nodes of the matched pattern.

\begin{Theorem}\label{thm:typesafe}
For the same pattern junction $\syn{J}$,
let $(\overline{r},\overline{q})$ be the result of function $\mb$ in \autoref{tab:st},
$\overline{r:s}$ the type assignment obtained by ({\sc Case}) in \autoref{tab:tc}.
We have $\overline{\atype(q) \subty s}$.
\end{Theorem}
\begin{proof}
We study the type $t$ of the node $q$ associated to a reference $r$ in function $\mb$ by case analysis.
In the case of ({\sc Match-One}), $t$ comes from the type assignment of pattern graph $P$.
By the definition of $\gmat[f]$, $t \subty s$, where $s$ is the type of node $R_1(r)$ in pattern graph $P_1$, such that $q = f(R_1(r))$.
In fact, in ({\sc Match-Some}), $P_1$ and $R_1$ are obtained from the type-checking rules.
According to ({\sc Pattern}), ({\sc Graph}) and ({\sc Node}), $s$ is exactly the type assigned to $r$ in the type-checking.

Similarly, in the case of ({\sc Match-Conj}), $P$ matches all $\overline{P_1}$.
The type $t$ of $q$ is a subtype of all the types $\overline{t_k}$ of the corresponding nodes in each of $\overline{P_1}$.
By \autoref{algo:conj-dfs-cons}, if there exists $\inf(\overline{t_k})$, it is assigned to $r$ in ({\sc Conj}),
for $\inf(\overline{t_k})$ is a subtype of each $\overline{t_k}$, and $t$ must be a subtype of the infimum by definition.
It there exists no infimum, $\NoType$ is assigned to the node in the union,
thus the $\overline{\min(\NoType,t_k)}$ in ({\sc Conj}) assigns $r$ with one of $\overline{t_k}$, of which $t$ is a subtype.

In the case of ({\sc Match-Disj}), $P$ matches only one graph $P'$ of $\overline{P_1}$.
The node $q$ bound to $r$ is non-deterministically chosen from nodes $\overline{f(u)}$ in $P$,
matching those nodes $\overline{u}$ in $P'$ that are finally mapped to the same node $\ov{u}$ in the intersection of $\overline{P_1}$.
The type $t$ of $q$ is a subtype of the corresponding node in $P'$ mapped to $\ov{u}$.
By \autoref{algo:disj-dfs-cons} and ({\sc Disj}), $r$ is assigned with the supremum type of all the node in $\overline{P_1}$ mapped to $\ov{u}$,
or $\NoType$ if the supremum does not exist.
Obviously, in either case, $t$ is a subtype of the type assigned to $r$.
\end{proof}

\endinput
